\documentclass[11pt,a4paper]{article}
\usepackage[T1]{fontenc}
\usepackage{lmodern,amsmath,amssymb,booktabs,microtype}
\usepackage[margin=24mm]{geometry}
\usepackage[colorlinks=true,urlcolor=blue,linkcolor=blue]{hyperref}
\newcommand{\EE}{\mathbb E}
\DeclareMathOperator{\tr}{tr}
\title{Parity-weighted rank deficits in the memory core\\Exact signed moments and a double-scaled Gram matrix}
\author{Alan Day \and Chengheng\thanks{AI research assistant (OpenAI).}}
\date{Local research note --- 6 November 2026}
\begin{document}\maketitle
\begin{abstract}
The output-rank deficit in the preceding memory formula is governed
by a symmetric-product selection rule. We derive its first two
parity-weighted moments from the existing Racah difference equation,
without assuming a complete Poisson law. This gives the limiting
Gram matrix of the energy-square, centered input-deficit, and
output-deficit operators. The output direction remains nonzero
after the energy-square projection and is orthogonal to the input
direction; its comparison with the force operator $A$ remains open.
\end{abstract}

\section{The finite half-filled rank problem}
Take $N=2m$, standard CAR, and the maximally mixed charge-$m$ state
with normalized trace $\tr_m$. Set $D=\binom Nm$,
$d_r=\binom Nr^2-\binom N{r-1}^2$, and $x_r=r(N+1-r)$.
The neutral Casimir has eigenvalue $2x_r$ on rank $r$.
For independent isotropic Hermitian operators of ranks $a,b$ with
$\EE U_a^2=\EE V_b^2=I$, use the previously derived positive law
\begin{equation}
 \pi_{ab}(r)=\EE\|P_r(U_aV_b)\|_m^2,\qquad
 \sum_r\pi_{ab}(r)=1,\qquad r\le a+b.
\end{equation}
The pure-rank covariance eigenvalues are
\begin{equation}
 \chi_a(b)={}_4F_3\!\left[
 \begin{matrix}-a,a-N-1,-b,b-N-1\\-m,-m,-N-1\end{matrix};1\right].
\end{equation}
They are the known Racah specialization~\cite{KS}. Let
\begin{equation}
 S=a+b-r,\qquad X=x_{a+b}-x_r,\qquad f=\chi_a(b).
\end{equation}
Then the exact signed moments, with expectation under $\pi_{ab}$, are
\begin{align}
 \EE_\pi(-1)^S&=f,\\
 \EE_\pi[(-1)^S X]&=-2abf,\\
 \EE_\pi[(-1)^S X^2]
   &=(4a^2b^2-2ax_b)f\notag\\
   &\quad+a^2(m+1-a)[\chi_{a-1}(b)-f].
 \label{eq:second}
\end{align}
For $a=0$ the last term is defined as zero. The formulas include
all output ranks and keep the sign; signed moments need not be positive.

\newpage
\section{A proof from the Racah transform}
At half filling the linear particle-hole anti-automorphism
$\tau(X)=\star X^{\mathsf T}\star^{-1}$ acts on rank $r$ by $(-1)^r$.
Let $T_\mu$ be a full Hermitian matrix basis orthonormal for normalized
trace. Completeness gives
\begin{equation}
 \sum_\mu(T_\mu)_{ij}(T_\mu)_{kl}=D\delta_{il}\delta_{jk},\qquad
 \sum_\mu\tau(T_\mu)XT_\mu=(-1)^mD\tau(X).
\end{equation}
The second identity uses $\star^{\mathsf T}=(-1)^m\star$.
Equivalently, for normalized pure-rank covariance maps $\Phi_r$,
\begin{equation}
 \sum_r(-1)^r d_r\Phi_r=(-1)^mD\tau,
 \qquad
 \sum_r(-1)^r d_r\chi_r(j)=(-1)^{m+j}D.
 \label{eq:transform}
\end{equation}

Specialize the existing difference equation~\cite[Eq.~(1.2.5)]{KS}
with $\alpha=\beta=-m-1$, $\gamma=-N-2$, $\delta=0$.
Define
\begin{align}
 (Lf)(j)&=B_j[f(j)-f(j+1)]+D_j[f(j)-f(j-1)],\\
 B_j&=\frac{(j-m)(j-N-1)^2}{2(2j-N-1)},&
 D_j&=\frac{j^2(j-m-1)}{2(2j-N-1)}.
\end{align}
The endpoints have $B_m=D_0=0$, and
\begin{equation}
 L\chi_a=x_a\chi_a,\quad
 2(B_j+D_j)=K-x_j,\quad Lx_j=N(x_j-K/2),
 \qquad K=\frac{N(N+1)}2.
\end{equation}
Writing $\varepsilon_j=(-1)^j$, conjugation gives
\begin{equation}
 \varepsilon L\varepsilon=K-x-L.
\end{equation}
Insert Eq.~\eqref{eq:transform} into the established Racah triple sum
for $\pi$. Multiplication by $x_r$ becomes $L$ on the transformed
index. If $\mathsf X$ denotes multiplication by $x_a$ and
$\sigma_k=\EE_\pi[(-1)^Sx_r^k]$, the first two operations give
\begin{align}
 \sigma_1&=(\mathsf X+L_a)\chi_a(b)=(x_a+x_b)f,\\
 \sigma_2&=\left[(\mathsf X+L_a)^2
             +N(K/2-\mathsf X-L_a)\right]\chi_a(b).
\end{align}
Use the recurrence once more to eliminate $f(a+1)$.
In $\EE(-1)^SX^2$, all order-$N^3$ terms cancel, leaving
Eq.~\eqref{eq:second}. This stable expression needs only
$\chi_a(b)$ and $\chi_{a-1}(b)$, rather than a full output-rank sum.

\newpage
\section{Signed and even-weighted double-scaled moments}
Let $p^2/N\to\lambda>0$, $a/p\to\alpha$, $b/p\to\beta$,
and write $\theta=2\lambda\alpha\beta$.
The terminating Racah series is dominated uniformly for $a,b\le p$
by $K_\lambda^k/k!$. Each fixed summand tends to
$(-4\lambda\alpha\beta)^k/k!$. Hence
\begin{equation}
 \chi_a(b)\longrightarrow e^{-2\theta},\qquad
 \chi_{a-1}(b)-\chi_a(b)\longrightarrow0.
\end{equation}
The prefactor $a^2(m+1-a)/N^2$ in Eq.~\eqref{eq:second} is bounded.
The three exact formulas consequently imply
\begin{equation}
 \left(\EE(-1)^S,\ \EE(-1)^S\frac XN,
                   \ \EE(-1)^S\frac{X^2}{N^2}\right)
 \longrightarrow e^{-2\theta}(1,-\theta,\theta^2-\theta).
\end{equation}
For $a+b\le2p$, the exact identity
\begin{equation}
 X=S[N+1-2(a+b)+S]
\end{equation}
gives $X/N=S[1+O(p/N)]$ uniformly. The previous positive moment
bounds therefore transfer the first two signed limits from $X/N$
to $S$. With the previous unsigned moments, this yields
\begin{align}
 \EE[1+(-1)^S]&\to1+e^{-2\theta},\\
 \EE[S(1+(-1)^S)]&\to\theta(1-e^{-2\theta}),\\
 \EE[S^2(1+(-1)^S)]&\to
             \theta^2(1+e^{-2\theta})+\theta(1-e^{-2\theta}).
 \label{eq:even}
\end{align}
These determine the Gram matrix needed below. They do not establish
a complete Poisson distribution.

The actual symmetric product selects precisely these even channels:
\begin{align}
 P_r(V_bU_a)&=(-1)^S P_r(U_aV_b),\\
 \EE\|P_r\{U_a,V_b\}\|_m^2
   &=4\mathbf1_{\{S\ {\rm even}\}}\pi_{ab}(r).
\end{align}
The first equation follows by applying $\tau$ to the ordered product.
It is an identity for the multiplication map, not an assertion that
all product multiplicities are absent. In particular, the $S=1$
channel is removed. Ignoring this selection rule would give the
wrong small-$\lambda$ strength for the output-deficit operator.

\newpage
\section{The same-coupling quadratic Gram matrix}
Return to the full overlapping complex Gaussian $p$-body Hamiltonian
$H=\sum_a h_a$, $\EE J_{IJ}J_{KL}=w\delta_{IL}\delta_{JK}$,
$\hbar=1$, at exact half filling. Write
$\EE h_a^2=C_aI$, $C=\sum_a C_a$, $\omega_a=C_a/C$,
and $\bar s_N=\sum_a(p-a)\omega_a$. Define
\begin{align}
 B_N&=W_2H^2=H^2-CI,\\
 L_N&=W_2\{H,D_N-\bar s_NH\},\qquad D_N=\sum_a(p-a)h_a,\\
 \mathcal S_N&=\sum_{a\le b,r}(a+b-r)W_2Z^+_{ab;r}.
\end{align}
Here $W_2$ means ordinary Gaussian coupling chaos, and $Z^+$ uses
an anticommutator for $a<b$ and a square for $a=b$.

For any two real symmetric rank filters $f,g$, orthogonality of
input pairs and output ranks, followed by Gaussian decoupling on
the diagonal pair, gives the exact identity
\begin{equation}
 \left\langle\sum fW_2Z^+,\sum gW_2Z^+\right\rangle
 =C^2\sum_{a,b}\omega_a\omega_b
       \sum_r[1+(-1)^S]\pi_{ab}(r)f_{ab;r}g_{ab;r}.
 \label{eq:gram}
\end{equation}
All scalar products are annealed normalized HS products.
Use $f=(1,u+v-2\bar s_N,S)$ for $(B_N,L_N,\mathcal S_N)$.
The known input deficits tend to independent Poisson variables of
mean $\lambda$ in this weighted sum, so $\bar s_N\to\lambda$.
For bounded input deficits Eq.~\eqref{eq:even} has
$\theta=2\lambda$ and $q=e^{-4\lambda}$. The input factorial tail
and output third-moment bound justify convergence of every Gram
entry, including the mixed entries. Thus
\begin{equation}
 \frac{G}{C^2}\longrightarrow
 \begin{pmatrix}
 1+q&0&\theta(1-q)\\
 0&2\lambda(1+q)&0\\
 \theta(1-q)&0&\theta^2(1+q)+\theta(1-q)
 \end{pmatrix}.
 \label{eq:matrix}
\end{equation}

Subtract the energy-square projection,
$\mathcal S_\perp=\mathcal S_N-\mu B_N$,
$\mu=\theta(1-q)/(1+q)$. In the physical normalization
$C\to J_s^2/\lambda$,
\begin{equation}
 \|\mathcal S_\perp\|_2^2\longrightarrow
 J_s^4\left[\frac{2(1-q)}\lambda+\frac{16q}{1+q}\right]>0,
 \qquad \langle\mathcal S_\perp,L_N\rangle\to0.
\end{equation}
After this fixed-$\lambda$ limit, the norm tends to $16J_s^4$ as
$\lambda\to0$. The old linear-filter bridge relates $L_N$ to
$J_HF$: $J_HF-4\lambda qC L_N\to0$ in $L^2$. The product extension
uses the Gaussian fourth-moment bound for the linear bridge error.
Hence $\langle\mathcal S_\perp,J_HF\rangle\to0$ as well.

\newpage
\section{Meaning, remaining work, and evidence}
The preceding memory-core result contains an output-deficit term
$8\lambda Cq^2\mathcal S_N$. We now know its equal-time norm and its
cross terms with $H^2$ and the input-deficit direction. The nonzero
projected norm shows that the output filter is not exhausted by
these two directions in the limiting HS space.

The cross term with the force observable
$A=\sum_i[H,c_i]^\dagger[H,c_i]$ is not determined by
Eq.~\eqref{eq:matrix}. It involves the coherent lower-charge product.
Consequently this calculation does not yet establish independence
from the full old span including $A$, a nonzero dynamical spectral
weight, or a new bulk particle. A projection on $H^2$ alone is not
the full projection onto conserved observables.

The next calculation should use the adjacent-charge covariance
and reduced-product recoupling to obtain the $A$ cross and weighted
tails. These are needed to replace its exact rank-dependent memory
coefficient by the fixed-deficit limit. Then the low-chaos pieces of
the fourth-chaos descendants can be restored in raw $Z_2$.
The full charge-window extension, finite-memory covariance, and
bulk space/mass dictionary remain part of the original objective.

\paragraph{Exact verification.}
For $N=2,4,6,8,10,12$ at half filling, all allowed input-rank pairs
give 417 signed-moment checks through second order. The same run
checks 139 complete Racah recurrences and the parity transform.
Independent CAR calculations use all symmetric Gaussian coordinate
pairs in three systems: 136 pairs for $(4,1,2)$ and 666 each for
$(4,2,2)$ and $(6,1,3)$. Their full $3\times3$ coefficient Grams
agree exactly with Eq.~\eqref{eq:gram}; these are not Monte Carlo
estimates.

\paragraph{Finite-size verification.}
Twenty-four cases at $\lambda=0.1,0.25,0.5$, $p=16,32,64,128$,
and input deficits $(0,0),(1,2)$ evaluate the exact finite moment
formulas at 80 and 120 digits. The measured finite variable is
$X/N$, not $S$. Actual finite-size deviations are retained.
At $\lambda=0.5,p=128,(u,v)=(1,2)$, the signed moment deviations
are approximately $(0.004323,-0.001067,-0.002125)$.
Decimal precision agreement is not interval certification.

The initial exact checker accidentally evaluated a negative integer
power of $-1$ in a zero-probability channel, causing conversion to
floating point. Replacing that expression by integer parity restored
exact arithmetic; the initial code is preserved. No formula or
tolerance was altered to obtain agreement.

\begin{thebibliography}{9}
\bibitem{KS} R. Koekoek and R. F. Swarttouw,
\emph{The Askey-scheme of hypergeometric orthogonal polynomials
and its $q$-analogue}, report 98-17 (1998), Sec.~1.2,
Eq.~(1.2.5),
\href{https://tudelft.roelofkoekoek.nl/documents/as98.pdf}{author-hosted report}.
\end{thebibliography}
\end{document}
